\scanpage{142}
\uline{Proposition} Pour que la suite $(H_i)$ d'él. hermitiens $\in L^p(\mathscr H)$ converge vers $H_0\in L^p(\mathscr H)$, il faut et il suffit que

\[
\text{1°)}\quad\|H_i\|_p\to\|H_0\|_p\qquad
\text{(ou seulement }\limsup_{i\to\infty}\|H_i\|_p\leq\unsure{\|H_0\|_p|\lvert H_0\|_p}
\]

2°) Pour \deleted{tout} $\lambda$ du spectre de $H_0$, (soit $\alpha>0$ tel qu'il \deleted{$[\lambda-\alpha,\lambda+\alpha]$ ne \unsure{contien|contienne}}\inserted{n'existe pas de} valeur spectrale $\mu\neq\lambda$ de $H_0$ telle que $|\lambda\deleted{\unsure{+}}-\mu|\leq\zeta$, \deleted{\unsure{élément spectral de $H$|l'ensemble spectral de $H$} \unsure{faut que}} \deleted{et telle} et tel que $\lambda-\alpha$ et $\lambda+\alpha$ n'appartiennent au spectre \deleted{\unsure{de} $H_{\unsure{0|i}}$}\inserted{$h_i$}\inserted{d'aucun $H_i$}) \deleted{\unsure{alors}} la composante spectrale de $H_i$ \deleted{\unsure{relative à}} dans $[\lambda-\alpha,\lambda+\alpha]$ est de dimension \deleted{\unsure{uniforme égale à}} qui finit par être égale à celle de $H_0$, et $h_i\to h_0$ (inutile de spécifier la topologie, en vertu d'une proposition préliminaire sur l'identité de \deleted{toutes} ces topologies)

\margin{\unsure{évidem\textsuperscript{t}|évident}\\ faux ds\\ le cas non\\ hermitien}

\uline{Corollaire 1} Les deux conditions suivantes sont néc. et suffis. \unsure{stes} pour la convergence de la suite $H_i$ vers $H_0$

\[
\text{1)}\quad\|H_i\|_p\to\|H_0\|_p\qquad\text{(resp$^{\unsure{\mathrm t|\mathrm d}}$ .\ .\ --)}
\]

\[
\text{2)}\quad H_i\to H_0\quad\text{dans }L^\infty(\mathscr H)
\]

\uline{Corollaire 2} $A\to |A|^r$ est application continue de $L^p$ dans $L^{\frac pr}$.

En effet, elle est composée de $A\to A^*A$, continue de $L^p$ dans $L^{\frac p2}$, avec
